As a consequence of $g$ being locally symmetric, the conformal group of the metric $g$ on $\rr^{8n}$ is contained in the group of homotheties,
\[H=\bigl\{ \phi\in \mathrm{Diff}\bigl(\R^{8n}\bigr)\mid \exists\ \lambda\in \rr\colon \psi^*g=\e^{2\lambda}g\bigr\}.\]
This is due to a result in \cite[Proposition 2.1]{CahenKerbrat77} that states that a conformal transformation between open sets in a semi-Riemannian manifold that is not conformally flat but has parallel Weyl tensor is a homothety. Since $g$ is locally symmetric it must also have parallel Weyl tensor, so that the result follows. Hence, we have the following.

\begin{Proposition}
The conformal group of $\bigl(\R^{8n},g\bigr)$, where $g$ is the metric in \eqref{metric} defined by the K\"{a}hler potential \eqref{potential} is equal to
\[\R\ltimes \Iso\bigl(\R^{8n},g\bigr),
\]
where the $\R$-factor is spanned by a one-parameter group of homotheties $\{\phi_s\mid s\in \R\}$.
Moreover, $g$~admits an abelian $2$-real-parameter family of homothetic transformations induced by the following diagonal linear maps of $\C^{4n}$,
\belabel{chomoth}\phi_{a,b}:=
%\begin{pmatrix} \e^{s+t}\1_{n} &0&0&0\\ 0& \e^{s-t} \1_{n}&0&0 \\
%0&0& \e^{3s-t} \1_{n}&0\\
%0&0&0&\e^{3s+t}\1_{n}
%\end{pmatrix} =
\begin{pmatrix} \e^{a}\1_{n} &0&0&0\\ 0& \e^{b} \1_{n}&0&0 \\
0&0& \e^{a+2b} \1_{n}&0\\
0&0&0&\e^{2a+b}\1_{n}
\end{pmatrix},\qquad a,b\in \R,\end{equation}
that fix the origin $($as any linear map$)$ and satisfy $\phi_{a,b}^*g=\e^{2(a+b)}g$.

\end{Proposition}
\begin{proof}
We begin by showing that the metric $g$ admits a homothety that is not an isometry.
Indeed, for each $s\in \R$, the linear diagonal transformation
\[\phi_s:=\begin{pmatrix} \e^{s}\1_{2n} &0\\ 0& \e^{3s}\1_{2n}\end{pmatrix}\]
of $\C^{4n}$ induces a homothety of $g$ with $\phi_s^*g=\e^{4s}g$.
Composing the homothety $\phi_s$ with the isometry induced by
\[\psi_t:=\begin{pmatrix} \e^{t}\1_{n} &0&0&0\\ 0& \e^{-t} \1_{n}&0&0 \\
0&0& \e^{-t} \1_{n}&0\\
0&0&0&\e^{t}\1_{n}
\end{pmatrix},\]
we obtain the {two-parameter} family of homotheties in the proposition as
\[\phi_{a,b}= \phi_{\tfrac{a+b}{2}} \circ \psi_{\tfrac{a-b}{2}}. \]
This completes the proof.
\end{proof}

We will now construct a compact conformal manifold with essential conformal transformations.

\btheo\label{main:thm}
Let $\tM=\rr^{8n}\setminus\{0\}$ and $g$ the K\"ahler metric on $\tM$ defined by the potential $f$ in~\eqref{potential} for a given non-zero symmetric matrix $(\sigma_{ij})$.
Let $a,b>0$ be fixed positive real numbers and $\Gamma_{a,b} $ be the cyclic group of {holomorphic} homotheties of $(\tM,g)$ generated by $\phi_{a,b}$ in~\eqref{chomoth}. Then
$M_{a,b}=\tM/\Gamma_{a,b}$ is a compact manifold diffeomorphic to $\SS^1\times \SS^{8n-1}$ and carries a $($non-flat$)$ conformal structure of signature $(4n,4n)$ that is locally conformal to a Ricci-flat K\"ahler metric and has essential conformal transformations.
$($The induced integrable complex structure on the quotient manifold preserves the induced conformal structure.$)$
\etheo
\bprf

Since $a$ and $b$ are positive, $M_{a,b}=\tM/\Gamma_{a,b}$ is a compact manifold that is diffeomorphic to $\SS^1\times \SS^{8n-1}$.
In fact, a fundamental domain diffeomorphic to $(0,1)\times \SS^{8n-1}$ is given by
\[ \mathcal D = \Biggl\{ (x_1,x_2,x_3,x_4) \in \mathbb{R}^{8n} = \bigl(\mathbb{R}^{2n}\bigr)^4 \mid 1< \sum_{i=1}^4 \|x_i\|^2 \  \mbox{and}\
\sum_{i=1}^4 \frac{1}{\lambda_i^2} \|x_i\|^2<1\Biggr\},\]
 where $\lambda_1 = \e^a$, $\lambda_2 = \e^b$, $\lambda_3 = \e^{a+2b}$, $\lambda_4 = \e^{2a+b}$, and $\|\, \cdot\,\|$ is the Euclidean norm on $\R^{2n}$.
 Note that this is the region enclosed by the standard sphere and an ellipsoid. Identifying the two
 boundary components of $\mathcal D$ (i.e., the sphere and the ellipsoid) via the diffeomorphism induced by $\phi_{a,b}$ yields a diffeomorphism $M_{a,b} \cong \SS^1\times \SS^{8n-1}$.

{As} $\Gamma_{a,b}$ is a group of homotheties, the metric $g$ defines a conformal structure $\mathbf{c}$ on {$M_{a,b}$.} This conformal structure is locally conformally Ricci-flat K\"ahler and not conformally flat.

For $t\in \rr$, consider the homotheties
\[\phi_{0,t}=
\begin{pmatrix} \1_{n} &0&0&0\\ 0& \e^{t} \1_{n}&0&0 \\
0&0& \e^{2t} \1_{n}&0\\
0&0&0&\e^{t}\1_{n}
\end{pmatrix},
\]
of $\bigl(\tM,g\bigr)$ with fixed point set $\bigl\{\bigl(z^1,\ldots , z^n,0,\ldots, 0\bigr)\mid z^i \in \C^n\setminus \{0\}\bigr\}$. Since the $\phi_{0,t}$ have fixed points, they are essential on $\tM=\rr^{8n}\setminus\{0\}$, see \cite[Proposition 2.5]{LeistnerTeisseire22}.
In addition, for each~${t\in \rr}$, these homotheties commute with $\phi_{a,b}$ and hence descend to a conformal map $\psi_t$ on the quotient~$M_{a,b}$ and are essential on the quotient. Indeed, if there was a metric in the conformal class on~$M_{a,b}$ for which $\psi_t$ {was} an isometry, then this metric would lift to $\tM$ and have~$\phi_{0,t}$ as an isometry. This is a contradiction as the $\phi_{0,t}$ are essential on $\tM$.
\end{proof}
